\section{Background}

\subsection{Feedback Quality in Medical Education}
Effective feedback in clinical education should be specific, behavior-oriented,
objective, and actionable. It should identify both strengths and areas for
improvement while avoiding vague praise or personal judgment. In OSCE examiner
training, these requirements are especially important because examiners must
provide concise, educationally useful feedback based on observed student
performance.

The distinction between feedback, encouragement, and evaluation is important for
this project. Encouragement may support motivation, but it often lacks
information value. Evaluation can summarize performance, but may remain too
summative to guide improvement. Formative feedback should instead describe a
performance gap and help the learner close it. In an OSCE setting, this means
that useful examiner feedback should name what the student did, connect the
comment to the station context, use descriptive rather than personal language,
and provide a concrete strategy for the next attempt.

These requirements become harder under the practical constraint of a two-minute
feedback window. The examiner must select the most relevant points, formulate a
strength, identify an improvement area, and give an actionable next step without
becoming either too vague or too lengthy. This makes feedback delivery a
trainable performance skill in its own right. A scalable training tool should
therefore allow examiners to rehearse feedback repeatedly and receive immediate,
criterion-based guidance.

This educational view motivates the six-criterion rubric used in Examiner
Coach. The rubric does not attempt to model all aspects of clinical competence.
It models the quality of the examiner's feedback utterance as a training object.
This separation is essential because the same student performance could receive
feedback of very different quality depending on the examiner's wording.

\subsection{Retrieval-Augmented Generation and Modular RAG}
Retrieval-augmented generation combines large language models with external
knowledge retrieval. Instead of relying only on model-internal knowledge, a RAG
system retrieves relevant evidence and provides it as context for generation.
The Modular RAG framework describes RAG systems as configurable combinations of
modules such as indexing, pre-retrieval, retrieval, post-retrieval, generation,
and orchestration \cite{gao2024modularrag}. Examiner Coach follows this idea by
separating its knowledge preparation, transcript processing, retrieval, scoring,
and coaching components.

The Modular RAG perspective is useful because Examiner Coach is not simply a
single prompt around a transcript. It is a pipeline with several points at which
behavior can be changed and evaluated independently: documents can be chunked
differently, retrieval can use direct transcript queries or HyDE-style
hypothetical-passage retrieval, retrieved chunks can be filtered or reranked, and generation can be
constrained by a structured JSON contract. This modularity improves
inspectability. When an evaluation seems poorly calibrated, the system can be
debugged at the level of retrieval evidence, prompt construction, parsing, or
rubric design rather than being treated as an opaque LLM response.

\subsection{KISSKI Context for Sensitive Health Applications}
The project also fits into the broader direction of KISSKI, the AI service
centre for sensitive and critical infrastructures. KISSKI's central goal is to
research AI methods and provide them through a highly available AI service
centre for sensitive infrastructures, with a particular focus on socially
relevant fields such as health and energy \cite{kisski2026mission}. Examiner
Coach takes a step in this direction by applying institutionally provided AI
services to a concrete health-education use case: OSCE examiner feedback
training. The project is not only a consumer of an LLM endpoint; it explores how
such services can be integrated into a structured, privacy-conscious workflow
for medical education.

The modular perspective also supports the privacy-oriented requirements of the
project. Instead of sending examiner input to a general commercial chatbot, the
system is configured around institutional KISSKI/SAIA services and a local
knowledge base. This does not remove all privacy concerns, but it makes the data
flow more explicit and gives the project a clearer path toward institutional
review.
